\documentclass[12pt,a4paper]{book} % article, report, book.

% Preámbulo: paquetes, comandos, entornos, estilos y título de página.
\usepackage{anyfontsize}
\usepackage[utf8]{inputenc} %Codificación en utf8%
\usepackage[spanish]{babel} %Hyphenation (Guionado) en español%
\usepackage[T1]{fontenc} %Codificación de fuente%

\usepackage{graphics,graphicx} %paquetes gráficos estándar%
\usepackage{color}

\setlength{\textwidth}{15cm}    % Anchura neta del texto (globalmente).
\setlength{\topmargin}{-2cm}     % top margin (Global).
\setlength{\textheight}{26cm}

\newcommand{\pfctitlename}{Guiones de prácticas sobre la plataforma Raspberry Pi}
\newcommand{\pfcanno}{2015}

\begin{document}

\begin{titlepage}
  \begin{center}
    \large
    \scshape
    \bfseries
    Universidad de Málaga \\
    Escuela Técnica Superior de Ingeniería Informática
  \end{center}

  \smallskip

  \begin{center}
    \includegraphics[width=4cm,keepaspectratio]{graphs/EscudoETSIIcolor.pdf}

    \bfseries
    \color[rgb]{0,0.18,0.39}
    \fontsize{15}{15}
    \selectfont E.T.S. INGENIERÍA\\
    \fontsize{20}{20}
    \selectfont INFORMÁTICA
  \end{center}

  \begin{center}
    \scshape
    \large
    \bfseries
    Ingeniero Técnico en Informática de Sistemas
  \end{center}

  \smallskip

  \begin{center}
    \huge
    \sffamily
    \bfseries
    \pfctitlename
  \end{center}

  \begin{center}
    \huge
    \sffamily
    Instructions for hands-on sessions with the Raspberry Pi platform
  \end{center}

  \bigskip

  \begin{center}
    \huge
    \sffamily
    Realizado por\\
    ANTONIO JOSÉ VILLENA GODOY
  \end{center}

  \begin{center}
    \huge
    \sffamily
    Dirigido por\\
    RAFAEL ASENJO PLAZA\\
    FRANCISCO J. CORBERA PEÑA
  \end{center}

  \begin{center}
    \huge
    \sffamily
    Departamento\\
    ARQUITECTURA DE COMPUTADORES
  \end{center}

  \vfill

  \hfill \large Málaga, Enero de \pfcanno
\end{titlepage}

\end{document}
